% --- 1. INFORMAZIONI GENERALI ---
\section{Informazioni Generali}
\begin{itemize}
	\item \textbf{Luogo/Piattaforma:} Server Discord\textsubscript{G}
	\item \textbf{Data:} 2026-04-28
	\item \textbf{Orario di inizio:} 14:30
	\item \textbf{Orario di fine:} 17:40
	\item \textbf{Responsabile\textsubscript{G}:} Filippo Idiometri
	\item \textbf{Scriba:} Sofia De Blasi
\end{itemize}

\vspace{0.5cm}
\textbf{Partecipanti presenti:}
\begin{table}[ht]
	\centering
	\renewcommand{\arraystretch}{1.2}
	\begin{tabularx}{\textwidth}{X l}
		\toprule
		\textbf{Nome e Cognome} & \textbf{Matricola} \\
		\midrule
		Sofia De Blasi & 2111014 \\
		Roberto Mariano Doroftei & 2111031 \\
		Filippo Idiometri & 2035617 \\
		Francesco Marcon & 2110990 \\
		Armando Moda Scarati & 2082864 \\
		Anton Ricardo Rupi & 2054304 \\
		\bottomrule
	\end{tabularx}
\end{table}

\vspace{0.5cm}
\textbf{Partecipanti assenti:}
\begin{table}[ht]
	\centering
	\renewcommand{\arraystretch}{1.2}
	\begin{tabularx}{\textwidth}{X l}
		\toprule
		\textbf{Nome e Cognome} & \textbf{Matricola} \\
		\midrule
		Nessuno & - \\
		\bottomrule
	\end{tabularx}
\end{table}

% --- 2. ORDINE DEL GIORNO ---
\section{Ordine del Giorno}
Gli argomenti previsti per la riunione odierna sono:
\begin{enumerate}
	\item Sprint\textsubscript{G} review e retrospective,
	\item Messa in pari sulla stesura dell'AdR e del PdP,
	\item Pianificazione dello sprint successivo e assegnazione dei task\textsubscript{G} e ruoli.
\end{enumerate}

% --- 3. RESOCONTO E DISCUSSIONE ---
\section{Resoconto della Riunione}

\subsection{Sprint review\textsubscript{G} e retrospective}
Il gruppo si è confrontato sull'andamento dello sprint appena terminato. Sono state segnalate
delle criticità riguardanti il Piano di Progetto\textsubscript{G} in particolare nell'aggiornamento delle tabelle 
che mostrano costi ed ore. Questo processo è, per ora, manuale e quindi sia dispendioso in termini di tempo 
che prono ad errori. Il gruppo ha quindi deciso di sviluppare un sistema di automazione di questi calcoli 
rendendo il lavoro del loro aggiornamento più semplice.
Sono emersi anche alcuni dubbi sulla definizione dei vari rischi e si è quindi deciso di aggiornare 
alcune parti della mitigazione dei rischi con riferimento ad essi.\\
Il gruppo ha inoltre deciso, per evitare errori di conteggio, uno standard sulle indicazioni delle tempistiche 
riguardanti le varie task e dei vari ruoli.\\
Infine il gruppo ha deciso di sviluppare il template\textsubscript{G} per il diario di bordo.

\subsection{Messa in pari sulla sulla stesura dell'AdR e del PdP}
Durante la riunione il gruppo ha effettuato un allineamento sullo stato di avanzamento dell'Analisi 
dei Requisiti (AdR) e del Piano di Progetto (PdP), con l'obiettivo di garantire una visione
condivisa del lavoro svolto e delle attività da proseguire nello sprint successivo.
L'allineamento ha permesso di portare alla luce alcuni dubbi riguardanti gli Use Case presenti nell' AdR 
che il gruppo ha deciso di usare come domande da porre al ricevimento con il professor Cardin.
A seconda dall'andamento del ricevimento il gruppo ha deciso di soffermarsi solamente sulla correzione e 
riscrittura dei vari casi d'uso\textsubscript{G} nel caso emergessero molte criticità e sull'avanzamento della stesura del 
documento altrimenti.\\
Riguardo al PdP invece non sono emerse criticità quindi il gruppo ha deciso di continuarne la sua stesura 
come previsto.

\subsection{Pianificazione dello sprint successivo e assegnazione dei task e ruoli}
Lo sprint è stato definito con durata dal 28 aprile al 5 maggio. I ruoli assegnati per questo sprint sono: 
Responsabile (Filippo Idiometri), Amministratore\textsubscript{G} (Anton R. Rupi), Verificatore\textsubscript{G} (Sofia De Blasi), 
Analisti (Armando M. Scarati, Roberto M. Doroftei, Francesco Marcon).
Come nei precedenti sprint, i ruoli di Programmatore\textsubscript{G} e Progettista\textsubscript{G} non sono ancora necessari e si è quindi 
deciso nuovamente di dedicare 3 membri al ruolo\textsubscript{G} dell'Analista\textsubscript{G}, in particolare Armando e Roberto riguardo all'
AdR e Francesco riguardo al PdP (stesura consuntivo sprint precedente).
% --- 4. DECISIONI PRESE ---
\section{Decisioni Prese}

\renewcommand{\arraystretch}{1.3}
\vspace{-0.3cm}
\noindent
\begin{longtable}{c p{12cm}}
	\toprule
	\textbf{Codice} & \textbf{Decisione} \\
	\midrule
	\endfirsthead

	\toprule
	\textbf{Codice} & \textbf{Decisione} \\
	\midrule
	\endhead

	\midrule
	\multicolumn{2}{r}{\textit{Continua nella pagina successiva}} \\
	\midrule
	\endfoot

	\bottomrule
	\endlastfoot

	\textbf{VI-11.1} & {Semplificazione del calcolo di tempi e costi per il PdP.} \\
	\textbf{VI-11.2} & {Standardizzazione delle tempistiche riguardanti le task.} \\
	\textbf{VI-11.3} & {Assegnazione dei rischi con codice e spiegazione quando necessari.}

\end{longtable}

% --- 5. AZIONI (TODO) ---
\section{Azioni da Svolgere (To-Do)}

\renewcommand{\arraystretch}{1.3}
\begin{longtable}{c c p{4.5cm} p{3cm} l}
	\toprule
	\textbf{Codice} & \textbf{Rif. Decisione} & \textbf{Descrizione Task} & \textbf{Assegnatario} & \textbf{Scadenza} \\
	\midrule
	\endfirsthead

	\toprule
	\textbf{Codice} & \textbf{Rif. Decisione} & \textbf{Descrizione Task} & \textbf{Assegnatario} & \textbf{Scadenza} \\
	\midrule
	\endhead

	\midrule
	\multicolumn{5}{r}{\textit{Continua nella pagina successiva}} \\
	\midrule
	\endfoot

	\bottomrule
	\endlastfoot

	% Issue\textsubscript{G} #137 - Aggiornamento AdR
	\ghissue{137} & - & Correzione degli errori, sviluppo di UC15, aggiornamento immagini & Armando M. Scarati, Roberto M. Doroftei & 2026-05-05 \\[4pt]
	
	% Issue #138 - Aggiornamento AdR
	\ghissue{138} & - & Aggiornamento delle Norme di Progetto & Sofia De Blasi & 2026-05-05 \\[4pt]
	
	% Issue #139 - Aggiornamento AdR
	\ghissue{139} & - & Stesura del verbale interno della settimana & Filippo Idiometri & 2026-05-05 \\[4pt]
	
	% Issue #140 - Aggiornamento AdR
	\ghissue{140} & - & Sviluppo del template del diario di bordo & Anton R. Rupi & 2026-05-05 \\[4pt]
	
	% Issue #141 - Aggiornamento AdR
	\ghissue{141} & VI-11.3 & Stesura dello sprint 3 nel PdP, aggiornamento rischi dello sprint 2 & Armando M. Scarati, Roberto M. Doroftei & 2026-05-05 \\[4pt]

\end{longtable}

% --- 6. PROSSIMA RIUNIONE ---
\section{Prossima Riunione}
\begin{itemize}
	\item \textbf{Data:} 2026-05-05
	\item \textbf{Ora:} 14:30
	\item \textbf{OdG preliminare\textsubscript{G}:}
	\begin{itemize}
		\item Revisione\textsubscript{G} e verifica dei task del quarto sprint;
		\item Avanzamento del Piano di Progetto;
		\item Avanzamento dell'Analisi dei Requisiti\textsubscript{G};
		\item Pianificazione del quinto sprint e assegnazione dei ruoli.
	\end{itemize}
\end{itemize}